\documentclass[a4paper]{article}

% Basic packages
\usepackage[fontset=fandol]{ctex}
\usepackage{amsmath, amssymb}
\usepackage{graphicx}
\usepackage[margin=2.5cm]{geometry}
\usepackage[most]{tcolorbox}
\usepackage{etoolbox}
\usepackage{listings}
\usepackage{hyperref}
\usepackage{booktabs}
\usepackage{subcaption}
\usepackage{float}
\usepackage{tikz}
\usetikzlibrary{positioning, arrows.meta}
\IfFileExists{pgfplots.sty}{
    \usepackage{pgfplots}
    \pgfplotsset{compat=1.18}
}{}

% Highlight boxes
\newtcolorbox{knowledgebox}[1]{
    enhanced, colback=blue!5!white, colframe=blue!75!black, colbacktitle=blue!75!black,
    coltitle=white, fonttitle=\bfseries, title=#1, attach boxed title to top left={yshift=-2mm, xshift=2mm},
    boxrule=1pt, sharp corners
}

\newtcolorbox{importantbox}[1]{
    enhanced, colback=yellow!10!white, colframe=yellow!80!black, colbacktitle=yellow!80!black,
    coltitle=black, fonttitle=\bfseries, title=#1, sharp corners
}

\newtcolorbox{warningbox}[1]{
    enhanced, colback=red!5!white, colframe=red!75!black, colbacktitle=red!75!black,
    coltitle=white, fonttitle=\bfseries, title=#1, sharp corners
}

% Listings
\lstset{
    language=Python,
    basicstyle=\ttfamily\small,
    keywordstyle=\color{blue},
    stringstyle=\color{red!60!black},
    commentstyle=\color{green!60!black},
    breaklines=true,
    frame=single,
    numbers=left,
    numberstyle=\tiny\color{gray},
    captionpos=b,
    extendedchars=false
}

% Front-page metadata
\newcommand{\notetitle}{CS224N Lecture 4: Dependency Parsing}
\newcommand{\noteauthors}{整理：AI Course Notes Contributors；授课：Chris Manning}
\newcommand{\notedate}{2024年1月}
\newcommand{\videochannel}{Stanford Online}
\newcommand{\videopublishdate}{2024}
\newcommand{\videoduration}{1:18:52}
\newcommand{\videourl}{https://www.youtube.com/watch?v=KVKvde-_MYc}
\newcommand{\videocoverpath}{}
\newcommand{\repourl}{https://github.com/hqhq1025/ai-course-notes}


% Footer with repo info
\usepackage{fancyhdr}
\pagestyle{fancy}
\fancyhf{}
\fancyfoot[C]{\thepage}
\fancyfoot[R]{\footnotesize\color{gray}\href{\repourl}{AI Course Notes}}
\renewcommand{\headrulewidth}{0pt}
\renewcommand{\footrulewidth}{0pt}

\begin{document}
\setlength{\emergencystretch}{2em}

\begin{titlepage}
\centering
{\Large 课程笔记\par}
\vspace{1.2cm}
{\huge\bfseries \notetitle\par}
\vspace{0.8cm}
{\large \noteauthors\par}
\vspace{0.3cm}
{\large \notedate\par}
\vspace{1.2cm}

\vspace{1.2cm}
\vfill
\begin{tcolorbox}[width=0.9\textwidth, colback=blue!3!white, colframe=blue!55!black, sharp corners]
\textbf{课程版本：} 2024 年中文讲义整理版\par
\textbf{笔记来源：} AI Course Notes Contributors\par
\textbf{本版处理：} 移除课件截图和对应图注，不包含 Stanford 原始课件、图片或视频。\par
\textbf{授权：} 笔记依 CC BY-NC-SA 4.0 分享；完整许可与改动记录见下一页。
\end{tcolorbox}
\end{titlepage}

\begin{center}
\fbox{\begin{minipage}{0.91\textwidth}
\small
\textbf{来源与许可}\\
本中文笔记由 AI Course Notes Contributors 整理，课程版本为 2024；源稿：\url{https://github.com/hqhq1025/ai-course-notes/tree/717e2df65c3d2eee593c171156a30d4f8f6812b9/cs224n/lecture04}。\\
笔记内容依 CC BY-NC-SA 4.0 分享：\url{https://creativecommons.org/licenses/by-nc-sa/4.0/}。\\
本副本的改动：移除 Stanford 原始课件截图与依赖这些截图的图注，移除课程视频封面/链接，并清理 TeX 源稿中的行尾空白后重新编译为可独立阅读的 PDF；同时加入本来源与许可说明。完整许可文本随本文件夹提供。Stanford 课件、课件图片和视频不包含在该授权内；请从对应年份 Stanford 课程页访问原始课件。
\end{minipage}}
\end{center}
\clearpage

\tableofcontents
\newpage

%% ========== 引言 ==========

\section{引言：为什么需要句法结构}

当我们阅读或听到一句话时，表面上是一个线性的词序列（或线性的声音流），但人类大脑会自动地将其组织为具有层次结构的意义单元——判断哪些词修饰哪些词、哪些词组成一个短语。自然语言处理（NLP）系统同样需要理解句子的内部结构，才能正确地解释语言。

\begin{importantbox}{本讲核心内容}
\begin{enumerate}
    \item \textbf{句法结构的两种表示}：短语结构（Constituency/Phrase Structure）与依存结构（Dependency Structure）
    \item \textbf{依存语法的基本概念}：中心词（head）与依存词（dependent）、依存关系类型
    \item \textbf{句法歧义}：介词短语附着歧义、并列结构作用域歧义等
    \item \textbf{依存句法分析方法}：基于转移的分析（Transition-based Parsing）、基于图的分析（Graph-based Parsing）
    \item \textbf{神经网络依存分析器}：用分布式表示替代离散特征，实现更高精度与更快速度
\end{enumerate}
\end{importantbox}

本讲内容直接关联课程作业 Assignment 2 的第二部分——使用 PyTorch 构建一个神经依存句法分析器。


% Raster figure omitted from this study copy.



\subsection{本章小结}

句法结构是理解自然语言的基础。人类在理解语言时自动地解析句子的层次结构，NLP 系统也需要具备这种能力。本讲将从语言学概念出发，逐步引入计算方法。

%% ========== 短语结构语法 ==========

\section{句法结构的两种视角}

\subsection{短语结构语法（Phrase Structure Grammar）}

语言学中表示句子结构的第一种方式是\textbf{短语结构}（Phrase Structure），由\textbf{上下文无关文法}（Context-Free Grammar, CFG）来形式化描述。

\begin{knowledgebox}{上下文无关文法（CFG）基本概念}
上下文无关文法通过一组\textbf{重写规则}描述语言的结构。每条规则将一个非终结符展开为终结符或非终结符的序列。例如：
\begin{itemize}
    \item $\text{NP} \rightarrow \text{Det} \ (\text{Adj})^* \ \text{N}$（名词短语 = 限定词 + 可选的形容词 + 名词）
    \item $\text{PP} \rightarrow \text{P} \ \text{NP}$（介词短语 = 介词 + 名词短语）
    \item $\text{VP} \rightarrow \text{V} \ \text{PP}$（动词短语 = 动词 + 介词短语）
    \item $\text{NP} \rightarrow \text{NP} \ \text{PP}$（名词短语也可以递归包含介词短语）
\end{itemize}
\end{knowledgebox}

以短语 ``the cuddly cat by the door'' 为例，其短语结构树自底向上构建：

\begin{itemize}
    \item ``the cuddly cat'' 构成一个名词短语（NP）
    \item ``the door'' 构成另一个名词短语
    \item ``by the door'' 构成一个介词短语（PP）
    \item 整体 ``the cuddly cat by the door'' 是一个更大的名词短语
\end{itemize}


% Raster figure omitted from this study copy.




% Raster figure omitted from this study copy.



\subsection{依存结构（Dependency Structure）}

本讲的重点是第二种表示方式——\textbf{依存结构}（Dependency Structure）。依存结构不关心短语的层次嵌套，而是直接刻画\textbf{词与词之间的修饰关系}。

\begin{importantbox}{依存结构的核心思想}
依存结构通过\textbf{有向弧}（directed arcs）连接句子中的词对，表示\textbf{中心词}（head）与\textbf{依存词}（dependent/modifier）之间的关系。例如在 ``look in the large crate in the kitchen by the door'' 中：
\begin{itemize}
    \item ``look'' 是整个句子的中心词（root）
    \item ``crate'' 是 ``in'' 的宾语，``in the large crate'' 修饰 ``look''
    \item ``large'' 和 ``the'' 修饰 ``crate''
    \item ``in the kitchen'' 和 ``by the door'' 也修饰 ``crate''
\end{itemize}
\end{importantbox}


% Raster figure omitted from this study copy.




% Raster figure omitted from this study copy.



\subsection{本章小结}

句法结构有两种主要表示方式：短语结构（基于上下文无关文法的树形嵌套）和依存结构（基于词间修饰关系的有向图）。两者可以相互转化，但本课程重点使用依存结构，因为它更直接地表达词与词之间的语义关系，且在 NLP 实践中更为流行。

%% ========== 句法歧义 ==========

\section{句法歧义：理解语言的核心挑战}

人类语言与编程语言的一个根本区别在于：\textbf{人类语言存在全局歧义（global ambiguity）}，而编程语言的歧义总能通过语法规则唯一消解。

\subsection{介词短语附着歧义（PP Attachment Ambiguity）}


% Raster figure omitted from this study copy.



以新闻标题 ``Scientists count whales from space'' 为例：

\begin{itemize}
    \item \textbf{解读 1}：``from space'' 修饰 ``count'' —— 科学家\textbf{从太空}数鲸鱼（正确解读）
    \item \textbf{解读 2}：``from space'' 修饰 ``whales'' —— 科学家数\textbf{来自太空的}鲸鱼
\end{itemize}

这种歧义源于英语中介词短语（PP）可以附着到不同的中心词上。由于 PP 在英语中极为常见，这类歧义也极为普遍。

\begin{warningbox}{PP 附着歧义的指数爆炸}
当句子中有多个连续的介词短语时，可能的解读数量按\textbf{Catalan 数}增长（指数级）。Manning 举了 Wall Street Journal 的例子：``The board approved its acquisition by Royal Trustco Ltd. of Toronto for \$27 a share at its monthly meeting'' 包含 4 个连续 PP，理论上有 \textbf{14 种}可能的依存结构（实际由上下文和世界知识消歧）。
\end{warningbox}


% Raster figure omitted from this study copy.



\subsection{并列结构作用域歧义（Coordination Scope Ambiguity）}


% Raster figure omitted from this study copy.



另一类常见歧义来自\textbf{并列结构}（coordination）的作用域。例如：

``Shuttle veteran and longtime NASA executive Fred Gregory appointed to board''

\begin{itemize}
    \item \textbf{解读 1}：Fred Gregory 是一个人，既是 shuttle veteran 又是 NASA executive
    \item \textbf{解读 2}：有两个人——一个 shuttle veteran 和 NASA executive Fred Gregory
\end{itemize}

类似的歧义还出现在 ``Doctor: No heart, cognitive issues'' 中——``No'' 的作用域是仅覆盖 ``heart'' 还是同时覆盖 ``heart'' 和 ``cognitive issues''。


% Raster figure omitted from this study copy.



\subsection{修饰语附着歧义}


% Raster figure omitted from this study copy.



不定式短语 ``to be used for Olympics beach volleyball'' 可以修饰 ``beach''（海滩将用于沙滩排球）或 ``body''（尸体将用于沙滩排球），产生截然不同的含义。

\begin{knowledgebox}{不同语言的歧义类型不同}
Manning 特别指出，不同语言的句法歧义类型并不完全相同。例如，中文的介词短语修饰动词时通常出现在动词之前，宾语在动词之后，因此不会出现英语中 PP 修饰动词 vs. 修饰宾语的歧义。但中文有自己独特的歧义类型。
\end{knowledgebox}

\subsection{歧义的实际影响}

Catalan 数列决定了 PP 附着歧义的组合爆炸速度：

\begin{table}[H]
\centering
\begin{tabular}{cc}
\toprule
\textbf{PP 数量} & \textbf{可能的解读数} \\
\midrule
1 & 2 \\
2 & 5 \\
3 & 14 \\
4 & 42 \\
5 & 132 \\
\bottomrule
\end{tabular}
\caption{PP 数量与可能解读数的关系（Catalan 数）}
\end{table}

尽管歧义数量指数增长，人类读者通常能毫不费力地选择正确解读——这依赖于世界知识、上下文和语用推理。NLP 系统需要从数据中学会类似的消歧能力。

\subsection{句法结构在 NLP 中的应用}


% Raster figure omitted from this study copy.



依存结构在实际 NLP 系统中有广泛应用。例如在生物信息学领域，可以通过依存分析从文本中提取蛋白质-蛋白质相互作用事实。动词 ``interacts'' 的主语和介词宾语分别指向两个交互的蛋白质，并列结构则帮助识别多重交互关系。

\subsection{本章小结}

句法歧义是自然语言的根本特征，包括 PP 附着歧义、并列结构作用域歧义和修饰语附着歧义等。歧义数量随句子复杂度指数增长，但人类能利用世界知识和上下文轻松消歧。NLP 系统需要从标注数据中学习这种消歧能力。理解并正确解析句法结构对信息抽取、语义理解等下游任务至关重要。

%% ========== 依存语法 ==========

\section{依存语法详解}

\subsection{依存语法的基本概念}


% Raster figure omitted from this study copy.



依存语法的核心元素：

\begin{importantbox}{依存语法的形式化定义}
\begin{itemize}
    \item \textbf{依存关系}（Dependency）：一个有向弧，从\textbf{中心词}（head/governor）指向\textbf{依存词}（dependent/modifier）
    \item \textbf{依存树}（Dependency Tree）：所有依存关系构成一棵有向树，有唯一的根节点（root）
    \item \textbf{关系标签}（Relation Labels）：每条弧上标注语法关系类型，如 nsubj（名词主语）、obj（宾语）、amod（形容词修饰语）、det（限定词）、obl（斜格）等
    \item \textbf{虚拟根节点}（ROOT）：通常在句子前添加一个虚拟的 ROOT 节点，其唯一依存词是句子的主要动词
\end{itemize}
\end{importantbox}


% Raster figure omitted from this study copy.



\subsection{依存语法的悠久历史}


% Raster figure omitted from this study copy.



\begin{knowledgebox}{依存语法的历史脉络}
依存语法是人类历史上最古老的句法形式化方法：
\begin{itemize}
    \item \textbf{Panini}（约公元前 4-8 世纪）：最早的语法学家，用依存关系描述梵语（Sanskrit）的语法结构。他的语法最初以\textbf{口述方式}传承数百年，令人惊叹
    \item \textbf{阿拉伯语法传统}（第一千年）：使用依存语法描述阿拉伯语
    \item \textbf{Reed-Kellogg 句子图解}：在美国教育中广泛使用的句子图解法，本质上是依存语法的一种变体
    \item \textbf{Lucien Tesni\`ere}（20世纪）：现代依存语法的奠基人，确立了从中心词指向依存词的箭头方向
    \item 相比之下，\textbf{短语结构语法/CFG} 是非常晚近的发明（1940s--1950s，由 Chomsky 等人系统化）
\end{itemize}
\end{knowledgebox}

Manning 特别提到了一个有趣的事实：计算机科学家通常从 Chomsky 层次结构了解 Chomsky，但 Chomsky 层次结构最初并非为了形式语言理论——而是为了论证人类语言不能用正则/有限状态文法来充分描述。

\subsection{箭头方向的约定}

在依存语法中，箭头方向存在两种约定：

\begin{itemize}
    \item \textbf{Head $\rightarrow$ Dependent}：Tesni\`ere 传统，也是本课程使用的方向
    \item \textbf{Dependent $\rightarrow$ Head}：另一些研究者使用的方向
\end{itemize}

\begin{warningbox}{注意箭头方向约定}
在阅读不同的依存语法文献时，必须首先确认所使用的箭头方向约定。不同的工具和标注方案可能采用不同的约定，混淆方向会导致对结构的完全误读。
\end{warningbox}

\subsection{依存分析的决策依据}

Manning 总结了依存分析器在做出依存决策时需要考虑的四类信息：

\begin{enumerate}
    \item \textbf{词汇亲和性（Bilexical Affinities）}：两个词之间是否存在合理的依存关系。例如 ``discussion $\leftarrow$ issues'' 是合理的，但 ``the $\leftarrow$ completed'' 则不合理
    \item \textbf{依存距离（Dependency Distance）}：大多数依存关系发生在\textbf{相邻或近距离}的词之间。长距离依存虽然存在，但较为少见
    \item \textbf{中间成分（Intervening Material）}：依存弧很少跨越动词或标点符号
    \item \textbf{中心词配价（Valency of Heads）}：中心词能接受多少个依存词。例如 ``broke'' 通常有一个主语（左侧）和一个可选的宾语（右侧），但不能接受多个宾语
\end{enumerate}


% Raster figure omitted from this study copy.



\subsection{投射性与非投射性（Projectivity）}


% Raster figure omitted from this study copy.



\begin{importantbox}{投射性与非投射性依存}
\begin{itemize}
    \item \textbf{投射性（Projective）}：依存弧不交叉，等价于上下文无关文法可生成的结构。大多数依存关系满足投射性
    \item \textbf{非投射性（Non-projective）}：依存弧存在交叉。例如 ``I'll give a talk tomorrow on neural networks'' 中，``on neural networks'' 修饰 ``talk''，``tomorrow'' 修饰 ``give''，两条弧交叉
\end{itemize}
疑问句中的 Wh-移位也常产生非投射依存：``Who did Bill buy the coffee from yesterday?'' 中 ``who'' 是 ``from'' 的宾语，但被移到了句首。
\end{importantbox}

本课程的作业中只处理投射性依存分析。

\subsection{本章小结}

依存语法通过中心词-依存词关系描述句子结构，具有悠久的历史传统。依存树具有唯一根节点和树形结构约束。分析器的决策依赖词汇亲和性、距离偏好、中间成分和配价约束。大多数依存关系是投射性的，但自然语言中也存在非投射性的情况。

%% ========== 树库 ==========

\section{标注数据与树库（Treebanks）}

\subsection{从规则到数据驱动}

早期的 NLP 系统依赖\textbf{手写规则}来解析句子结构。这种方法存在两个致命问题：

\begin{enumerate}
    \item \textbf{覆盖率不足}：人类语言具有无穷的创造性（如 Yoda 式语序重排、网络用语等），手写规则难以穷举
    \item \textbf{无法消歧}：纯规则系统可以生成所有合法的句法分析结果，但无法在多个合法分析之间选择最可能的那个
\end{enumerate}


% Raster figure omitted from this study copy.



为了解决这些问题，NLP 社区从 1980 年代末到 1990 年代开始大规模建设\textbf{标注语料库}（annotated corpora），即\textbf{树库}（Treebanks）——由人工标注了句法结构的句子集合。

\begin{knowledgebox}{树库为何重要}
树库的建设虽然缓慢而费力，但带来了两大革命性优势：
\begin{enumerate}
    \item \textbf{可复用的训练数据}：一次标注，多次使用。可以训练句法分析器、词性标注器、心理语言学模型等
    \item \textbf{系统化的评估方法}：提供了客观的基准测试。在树库出现之前（1950s--1980s），NLP 系统的评估方式是：运行程序，输入一句话，看看输出的分析结果是否``看起来合理''——没有任何系统化的定量评估
\end{enumerate}
\end{knowledgebox}

\subsection{Universal Dependencies（UD）}


% Raster figure omitted from this study copy.



Manning 介绍了他深度参与的 \textbf{Universal Dependencies}（UD）项目——一个统一标注方案的跨语言依存树库：

\begin{itemize}
    \item 覆盖超过 \textbf{100 种语言}
    \item 使用统一的依存关系标签体系
    \item 适用于跨语言研究、心理语言学实验等
    \item 课程作业 Assignment 2 使用的就是这类数据
\end{itemize}

\subsection{本章小结}

树库的建设推动了 NLP 从规则驱动转向数据驱动。Universal Dependencies 项目提供了覆盖 100 多种语言的标准化依存树库，是现代依存分析研究和应用的基础资源。树库同时提供了训练数据和评估基准。

%% ========== 基于转移的依存分析 ==========

\section{基于转移的依存分析（Transition-based Parsing）}

\subsection{依存分析的方法概览}

依存分析有多种算法范式：

\begin{enumerate}
    \item \textbf{动态规划方法}（Dynamic Programming）
    \item \textbf{基于图的方法}（Graph-based / Maximum Spanning Tree）
    \item \textbf{约束满足方法}（Constraint Satisfaction）
    \item \textbf{基于转移的方法}（Transition-based / Shift-Reduce）——本讲重点
\end{enumerate}

基于转移的方法在实践中最为流行，因为它具有\textbf{线性时间复杂度}，远优于上下文无关文法分析的立方时间复杂度。

\subsection{Arc-Standard 转移系统}

基于转移的依存分析使用两个核心数据结构和三种操作：

\begin{importantbox}{转移系统的组成}
\textbf{数据结构}：
\begin{itemize}
    \item \textbf{栈}（Stack）：存放正在处理的词，初始只包含 ROOT。\textbf{栈顶在右侧}
    \item \textbf{缓冲区}（Buffer）：存放尚未处理的词，初始包含句子的所有词。\textbf{缓冲区顶在左侧}
\end{itemize}

\textbf{三种操作}：
\begin{enumerate}
    \item \textbf{SHIFT}：将缓冲区顶部的词移到栈顶
    \item \textbf{LEFT-ARC}：栈顶第二个词成为栈顶词的依存词（左弧），弹出第二个词
    \item \textbf{RIGHT-ARC}：栈顶词成为栈顶第二个词的依存词（右弧），弹出栈顶词
\end{enumerate}

\textbf{终止条件}：缓冲区为空，栈中只剩 ROOT。
\end{importantbox}

\subsection{实例演示：分析 ``I ate fish''}


% Raster figure omitted from this study copy.



以句子 ``I ate fish'' 为例，分析过程如下：

\begin{table}[H]
\centering
\small
\begin{tabular}{clll}
\toprule
\textbf{步骤} & \textbf{栈} & \textbf{缓冲区} & \textbf{操作} \\
\midrule
0 & [ROOT] & [I, ate, fish] & 初始状态 \\
1 & [ROOT, I] & [ate, fish] & SHIFT \\
2 & [ROOT, I, ate] & [fish] & SHIFT \\
3 & [ROOT, ate] & [fish] & LEFT-ARC（I $\leftarrow$ ate） \\
4 & [ROOT, ate, fish] & [] & SHIFT \\
5 & [ROOT, ate] & [] & RIGHT-ARC（ate $\rightarrow$ fish） \\
6 & [ROOT] & [] & RIGHT-ARC（ROOT $\rightarrow$ ate） \\
\bottomrule
\end{tabular}
\caption{``I ate fish'' 的转移操作序列}
\end{table}


% Raster figure omitted from this study copy.



最终构建的依存关系集合为：
\begin{itemize}
    \item ate $\rightarrow$ I（主语）
    \item ate $\rightarrow$ fish（宾语）
    \item ROOT $\rightarrow$ ate（句子根节点）
\end{itemize}

\begin{warningbox}{操作选择决定了分析结果}
不同的操作序列会生成不同的依存结构。如果在第 3 步选择 RIGHT-ARC 而非 LEFT-ARC，就会让 ``I'' 成为中心词而 ``ate'' 成为依存词——这显然是错误的。因此，\textbf{如何在每一步选择正确的操作}是整个系统的核心问题。
\end{warningbox}

\subsection{Nivre 与贪心决策}


% Raster figure omitted from this study copy.



瑞典 NLP 研究者 \textbf{Joakim Nivre}（2003--2005）提出了用\textbf{机器学习}来预测每一步应采取的操作。核心洞察：

\begin{importantbox}{贪心转移分析的关键发现}
\begin{enumerate}
    \item 每步操作都可以视为一个\textbf{分类问题}：给定当前状态（栈、缓冲区内容），预测应执行 SHIFT、LEFT-ARC 还是 RIGHT-ARC
    \item 尽管使用\textbf{贪心策略}（每步只选最优操作，不做搜索），分析精度仍然很高
    \item 整个分析过程是\textbf{线性时间} $O(n)$：每个词最多被 SHIFT 一次、参与一次 ARC 操作
\end{enumerate}
\end{importantbox}

\subsection{传统特征工程方法}

Nivre 最初（2005 年）使用的是传统的\textbf{符号化特征}（indicator features）配合线性分类器（如逻辑回归或 SVM）：


% Raster figure omitted from this study copy.



典型特征包括：
\begin{itemize}
    \item 栈顶词是 ``good'' 且词性是形容词
    \item 栈顶词是 ``good'' 且栈中第二个词是 ``had''
    \item 缓冲区顶部词的词性是名词
\end{itemize}

这些特征的组合可以达到数百万甚至更多。

\begin{warningbox}{传统特征方法的三大问题}
\begin{enumerate}
    \item \textbf{稀疏性}（Sparsity）：大多数特征组合在训练数据中极少出现，可能仅出现 10 次/百万句
    \item \textbf{不完备性}（Incompleteness）：有些词和组合在训练数据中从未出现
    \item \textbf{计算开销}：运行时分析中，\textbf{95\% 以上的时间}花在计算特征上，而非机器学习决策本身
\end{enumerate}
\end{warningbox}

\subsection{本章小结}

基于转移的依存分析通过栈和缓冲区模拟分析过程，使用 SHIFT、LEFT-ARC、RIGHT-ARC 三种操作逐步构建依存树。Nivre 提出的贪心策略实现了线性时间复杂度。传统方法依赖大量手工特征，存在稀疏性、不完备性和高计算开销等问题。

%% ========== 神经依存分析 ==========

\section{神经网络依存分析器}

\subsection{从离散特征到分布式表示}


% Raster figure omitted from this study copy.



针对传统特征方法的三大问题，自然的解决方案是使用\textbf{分布式表示}（distributed representations）——即我们在前几讲中学习的\textbf{词向量}（word embeddings）。

\subsection{Chen \& Manning (2014) 神经依存分析器}


% Raster figure omitted from this study copy.



Danqi Chen（当时是 Manning 的博士生，也是 CS224N 的两届 Head TA）构建了第一个成功的\textbf{神经转移依存分析器}，其核心设计如下：

\begin{importantbox}{神经依存分析器的架构}
\textbf{输入表示}：对于当前分析状态，提取以下关键元素的分布式表示并拼接（concatenate）：
\begin{itemize}
    \item 栈顶和缓冲区顶部的\textbf{词向量}（word embeddings）
    \item 对应位置的\textbf{词性向量}（POS tag embeddings）
    \item 已构建的依存弧的\textbf{关系标签向量}（dependency label embeddings）
\end{itemize}

\textbf{网络结构}：
$$h = \text{ReLU}(W_1 \cdot [\text{concatenated inputs}] + b_1)$$
$$\text{output} = \text{softmax}(W_2 \cdot h + b_2)$$

\textbf{输出}：一个概率分布，覆盖所有可能的操作（SHIFT、各种标签的 LEFT-ARC、各种标签的 RIGHT-ARC）。
\end{importantbox}


% Raster figure omitted from this study copy.



\subsection{三种分布式表示的作用}

\begin{knowledgebox}{为什么需要三种嵌入}
\begin{itemize}
    \item \textbf{词向量}（Word Embeddings）：捕获词汇语义相似性。即使训练数据中没有见过某个词在特定位置出现，如果该词与已见过的词语义相近，系统也能做出合理预测
    \item \textbf{词性向量}（POS Tag Embeddings）：将细粒度的词性标签（如单数名词 NN vs. 复数名词 NNS）映射到连续空间，使相似词性获得相近表示
    \item \textbf{依存标签向量}（Dependency Label Embeddings）：捕获已构建的依存关系类型之间的相似性，为后续决策提供结构化上下文
\end{itemize}
\end{knowledgebox}

\subsection{非线性分类器的优势}

传统方法使用\textbf{线性分类器}（逻辑回归、SVM），而神经方法使用\textbf{非线性分类器}（含 ReLU 隐藏层的前馈网络）。这带来了两大优势：

\begin{enumerate}
    \item \textbf{自动学习特征组合}：隐藏层能自动学习输入特征的非线性组合，不再需要手工设计特征交叉
    \item \textbf{更强的表达能力}：非线性分类器能拟合更复杂的决策边界
\end{enumerate}

\subsection{实验结果}


% Raster figure omitted from this study copy.



实验结果令人瞩目：

\begin{table}[H]
\centering
\begin{tabular}{lccc}
\toprule
\textbf{分析器} & \textbf{UAS (\%)} & \textbf{LAS (\%)} & \textbf{速度} \\
\midrule
MaltParser（Nivre, 转移+符号特征） & 89.9 & -- & 快 \\
MSTParser（图方法+符号特征） & 91.5 & -- & 慢（$\sim$50x） \\
Chen \& Manning（转移+神经网络） & 92.0 & -- & 快 \\
\bottomrule
\end{tabular}
\caption{不同依存分析器的精度与速度对比}
\end{table}

\begin{importantbox}{神经方法的双重突破}
Chen \& Manning (2014) 的神经依存分析器同时实现了：
\begin{enumerate}
    \item \textbf{精度提升}：超越了传统的转移分析器（MaltParser），达到了图方法（MSTParser）的水平
    \item \textbf{速度保持}：保持了转移分析的线性时间复杂度，甚至比传统转移分析器更快（因为省去了大量特征计算开销）
\end{enumerate}
这打破了精度-速度的传统权衡。
\end{importantbox}

\subsection{本章小结}

Chen \& Manning (2014) 的神经依存分析器用分布式表示替代离散特征、用非线性网络替代线性分类器，在保持线性时间复杂度的同时，将分析精度提升到图方法的水平。三种嵌入（词、词性、依存标签）共同提供丰富的上下文信息。这项工作开创了神经句法分析的新时代。

%% ========== 评估方法 ==========

\section{依存分析的评估}

\subsection{UAS 与 LAS}


% Raster figure omitted from this study copy.



依存分析的评估非常直观——逐条检查系统预测的依存弧是否与标准答案一致：

\begin{importantbox}{两种评估指标}
\begin{itemize}
    \item \textbf{UAS（Unlabeled Attachment Score）}：只看依存弧的方向是否正确（即每个词的中心词是否预测对了），不考虑关系标签
    $$\text{UAS} = \frac{\text{正确预测中心词的词数}}{\text{总词数}}$$
    \item \textbf{LAS（Labeled Attachment Score）}：同时要求依存弧方向和关系标签都正确
    $$\text{LAS} = \frac{\text{中心词和标签都正确的词数}}{\text{总词数}}$$
\end{itemize}
LAS $\leq$ UAS，因为 LAS 的要求更严格。
\end{importantbox}

例如对句子 ``She saw the video lecture''（5 个词），如果系统有 4 个词的中心词预测正确（1 个错误），但只有 2 个词的标签也完全正确，则 UAS = 4/5 = 80\%，LAS = 2/5 = 40\%。

\subsection{本章小结}

UAS 和 LAS 是依存分析的标准评估指标。UAS 只评估依存弧方向的正确性，LAS 同时要求方向和标签都正确。这种基于树库的定量评估方法是 1990 年代以来 NLP 领域的重要进步。

%% ========== 后续发展 ==========

\section{后续发展：从 SyntaxNet 到图方法}

\subsection{Google SyntaxNet / Parsey McParseface}


% Raster figure omitted from this study copy.



2016 年，Google 在 Chen \& Manning 的基础上进行了大规模改进，推出了 \textbf{SyntaxNet}（别名 \textbf{Parsey McParseface}）：

\begin{itemize}
    \item \textbf{更深的网络}：使用更多层的神经网络和更大的向量维度
    \item \textbf{更优的超参数}：经过大规模调参
    \item \textbf{柱搜索}（Beam Search）：不再是纯贪心策略，而是同时维护多个候选分析路径，最终选择最优的一条
    \item UAS 从约 92\% 提升到约 \textbf{94.6\%}
\end{itemize}

Manning 提到，Google 将一个依存分析器做成了大新闻——在 Wired、VentureBeat 等科技媒体上广泛报道——这让他感到意外。起了一个搞笑的名字（Parsey McParseface）确实有助于获得媒体关注。

\subsection{神经图方法依存分析（Neural Graph-based Parsing）}


% Raster figure omitted from this study copy.



除了转移方法，另一种重要的依存分析范式是\textbf{图方法}（Graph-based Parsing）：

\begin{importantbox}{图方法的核心思想}
\begin{enumerate}
    \item 对句子中的\textbf{每对词}，计算一个依存弧的得分（共 $O(n^2)$ 对）
    \item 每个词需要选择一个中心词（或者是 ROOT），评估所有候选中心词的得分
    \item 使用\textbf{最小生成树}（Minimum Spanning Tree）算法从得分矩阵中找出得分最高的合法依存树（无环、连通、唯一根）
\end{enumerate}
\end{importantbox}


% Raster figure omitted from this study copy.



图方法的得分计算也可以使用神经网络。Manning 的团队将图方法与神经网络结合，构建了比 Parsey McParseface 更准确的分析器（UAS 再提升约 1\%）。


% Raster figure omitted from this study copy.



这个神经图方法分析器被集成到了 Stanford 的开源 NLP 工具包 \textbf{Stanza} 中。

\begin{knowledgebox}{转移方法 vs. 图方法}
\begin{itemize}
    \item \textbf{转移方法}：线性时间 $O(n)$，速度快，适合大规模处理；贪心决策可能导致错误传播
    \item \textbf{图方法}：$O(n^2)$ 或 $O(n^3)$ 时间，精度通常更高；考虑全局结构约束（树约束）
    \item 在传统和神经两个时代，图方法的精度都略高于转移方法
    \item 实践中两种方法都被广泛使用，取决于对速度和精度的需求
\end{itemize}
\end{knowledgebox}

\subsection{本章小结}

Google 的 SyntaxNet 通过更深网络和柱搜索在转移方法上取得了进一步提升。图方法通过评估所有词对的依存可能性并求解最优树，提供了另一种精度更高的分析范式。两种方法的神经网络版本都显著优于传统符号化方法。

%% ========== 总结与延伸 ==========

\section{总结与延伸}

\subsection{讲者的核心总结}


% Raster figure omitted from this study copy.



Chris Manning 在本讲中构建了从语言学概念到计算方法的完整链条：

\begin{enumerate}
    \item \textbf{句法结构是理解语言的基础}：人类语言不是简单的词序列，而是具有层次结构的系统
    \item \textbf{依存语法是高效的结构表示}：直接刻画词间修饰关系，历史悠久，在现代 NLP 中广泛应用
    \item \textbf{歧义是核心挑战}：人类语言充满全局歧义，消歧需要世界知识和统计信息
    \item \textbf{基于转移的分析是实用的解决方案}：线性时间复杂度，适合大规模处理
    \item \textbf{神经网络带来了范式转换}：分布式表示解决了稀疏性问题，非线性分类器提升了精度，同时降低了计算开销
\end{enumerate}

\subsection{全课知识图谱}

\begin{center}
\begin{tikzpicture}[node distance=0.9cm, every node/.style={draw, rounded corners, minimum width=4cm, minimum height=0.7cm, font=\small}]
\node (ling) {句法结构（CFG vs. 依存语法）};
\node (amb) [below=of ling] {句法歧义与消歧};
\node (tb) [below=of amb] {树库与数据驱动方法};
\node (trans) [below=of tb] {基于转移的分析};
\node (neural) [below=of trans] {神经依存分析器};
\node (graph) [below=of neural] {图方法与后续发展};
\draw[->, thick] (ling) -- (amb);
\draw[->, thick] (amb) -- (tb);
\draw[->, thick] (tb) -- (trans);
\draw[->, thick] (trans) -- (neural);
\draw[->, thick] (neural) -- (graph);
\node[draw=none, right=0.8cm of ling, text width=5.5cm, font=\scriptsize] {短语结构、依存结构、中心词与依存词};
\node[draw=none, right=0.8cm of amb, text width=5.5cm, font=\scriptsize] {PP 附着、并列作用域、Catalan 数};
\node[draw=none, right=0.8cm of tb, text width=5.5cm, font=\scriptsize] {Universal Dependencies、100+ 语言};
\node[draw=none, right=0.8cm of trans, text width=5.5cm, font=\scriptsize] {Stack + Buffer、SHIFT/LEFT/RIGHT-ARC};
\node[draw=none, right=0.8cm of neural, text width=5.5cm, font=\scriptsize] {词向量 + 词性向量 + 标签向量、ReLU};
\node[draw=none, right=0.8cm of graph, text width=5.5cm, font=\scriptsize] {SyntaxNet、MST、Stanza};
\end{tikzpicture}
\end{center}

\subsection{关键 Takeaways}

\begin{importantbox}{五条核心原则}
\begin{enumerate}
    \item \textbf{依存语法 > 短语语法}（在 NLP 实践中）：依存结构直接表达词间关系，标注效率高，跨语言适用性强
    \item \textbf{数据驱动 > 规则驱动}：树库提供的统计信息远比手写规则更能有效消歧
    \item \textbf{分布式表示 > 离散特征}：词向量解决了稀疏性和不完备性问题，同时降低了计算开销
    \item \textbf{非线性分类器 > 线性分类器}：神经网络的隐藏层能自动学习特征组合
    \item \textbf{Transformer 中的隐式句法}：后续课程将讨论，Transformer 语言模型虽然不显式构建依存树，但其注意力机制隐式地学习了类似的结构信息
\end{enumerate}
\end{importantbox}

\subsection{拓展阅读}

\begin{itemize}
    \item Nivre, J. (2005). Dependency grammar and dependency parsing. MSI report 5133(1).
    \item Chen, D. \& Manning, C.D. (2014). A Fast and Accurate Dependency Parser using Neural Networks. \textit{EMNLP 2014}. \url{https://nlp.stanford.edu/pubs/emnlp2014-depparser.pdf}
    \item Andor, D. et al. (2016). Globally Normalized Transition-Based Neural Networks. \textit{ACL 2016}.（Google SyntaxNet / Parsey McParseface）
    \item Dozat, T. \& Manning, C.D. (2017). Deep Biaffine Attention for Neural Dependency Parsing. \textit{ICLR 2017}.（神经图方法）
    \item Universal Dependencies 项目：\url{https://universaldependencies.org/}
    \item Stanford Stanza NLP 工具包：\url{https://stanfordnlp.github.io/stanza/}
    \item Jurafsky, D. \& Martin, J.H. \textit{Speech and Language Processing}, Chapter 18: Dependency Parsing. \url{https://web.stanford.edu/~jurafsky/slp3/}
\end{itemize}

\end{document}
